\subsection{Consumer-side feasibility study}
A registry graph assigns downstream reach through observed dependency edges. To check how much of that reach can be traced in application-owned artifacts, we drew 40 repositories without replacement from 199 distinct GitHub repositories in the Awesome-Selfhosted index tagged JavaScript, TypeScript or Nodejs. We froze the index revision (\code{2a8926e0}), deterministic draw and midnight UTC cutoff on November 20, 2026 before inspecting consumer evidence. Each repository was pinned to a default-branch commit at or before the cutoff. None of the 40 appeared among the 7,171 original endpoint keys or the 66 endpoints used to train the October cohort rankings. The frame represents curated public self-hosted software.

Of 40 repositories, 34 supplied a root manifest and 17 a root npm lockfile. We evaluated eight with sufficiently supported inputs. Across the full draw, six lacked a root manifest, ten declared another manager, seven additional cases lacked a root npm lockfile, eight had unsupported inputs, and one contained an ambiguous or external path. These categories are mutually exclusive. For npm lockfile versions 2 and 3 we traced mandatory production dependencies from the root importer and checked declared ranges against locked versions. We compared each complete path with edges in npm Arborist's offline virtual tree. Both parsers use npm SemVer. This is a software cross-check of committed metadata. No packages were installed and no authors supplied manual labels.

Seven of eight compared projects yielded paths. The two parsers agreed on all 1,147 observed installation locations across those projects. Counting distinct package/version pairs separately per consumer yielded 1,110 observations. Of those, 278 could be linked through a public npm tarball identity to exact package/version maintainer metadata captured in the October cohort. Six projects supplied at least one such link. That metadata describes the earlier registry observation and cannot identify who controlled publication at a later date.

We evaluated the two greedy coverage portfolios frozen on the October cohort against these consumer paths. At budget 100, the name-based and version-based methods selected exactly the same 100 accounts. Both portfolios had a documented path in the same five projects. Three compared projects had no documented hit, and 32 drawn projects were unevaluable. Identical account lists necessarily cover the same projects when evaluated against the same files. The pilot shows how many projects we could evaluate and where the evidence ran out.
